\documentclass[../main]{subfiles}
\begin{document}
\ifSubfilesClassLoaded{\setcounter{page}{36}}{}
\section{Необхідна праця та умови для комуністичної діяльності}
\label{sec:06_necessitatem_truda}

У комуністичному суспільстві життя людини не опосередковується приватною власністю та обміном товарами, отже, воно вільне від усього того, що потрібно для забезпечення умов існування приватної власності. Це дає суспільству значні переваги та зберігає його трудові ресурси. У «Ельберфельдських промовах» Енгельс навів низку прикладів, які наочно показують, скільки людської праці можна зберегти в суспільстві, де інтереси членів не суперечать один одному, а продукти діяльності не перетворюються на товари, які протистоять один одному на ринку.

Праця, необхідна для функціонування ринку, через який розподіляються суспільні ресурси, замінюється ефективнішим прямим розподілом. Це дає можливість не здійснювати дій, які є безглуздими з погляду організації виробництва в суспільному масштабі, і не витрачати на них час людей. Виготовляти пакування, яке не потрібне з погляду простого зберігання продуктів, здійснювати перевезення, пов’язані з продажами, а не з доставкою кінцевому споживачу, виробляти продукти, які потім виявляться зайвими на ринку і не будуть продані, і тим більше виробляти те, що завдає шкоди людині, але приносить величезні прибутки в суспільстві товарного виробництва (наркотики тощо), – все це не має жодного сенсу в суспільстві, де виробництво не має товарного характеру.

Ліквідація приватної форми присвоєння також дає значну економію життєвої людської праці. Це те саме, що й товарне виробництво, але розглянуте з точки зору праці, витраченої на доведення продукту до кінцевого споживання. Приклади з повсякденного життя та транспорту, про які ми вже писали, прямо вказують на можливість вивільнення великої кількості людської праці, яка зараз неефективно використовується в індивідуальному порядку.

Сучасний капіталізм вимагає, щоб зусилля людей були спрямовані на підтримку можливості виробництва додаткової вартості. Кількість такої праці постійно зростає. Виробництво додаткової вартості неминуче доповнилося виробництвом споживання та способом життя, спрямованим на споживання певних товарів. Вже це свідчить про зростаючі (і значно більші, ніж за часів Маркса та Енгельса) можливості для звільнення людей, які зможуть на засадах колективної власності займатися виробництвом того, що потрібно суспільству для його розвитку.

Суспільний час витрачається на підтримку існування тих обмежень, які приватна власність накладає на життя суспільства, та на усунення її наслідків, шкідливих для суспільства, зокрема й злочинів: поліція, регулярна армія, чиновники, різноманітні охоронці приватної власності. Усі ці трудові ресурси також можуть бути звільнені після подолання приватної власності на засоби виробництва та перенаправлені на виконання праці, в якій суспільство безпосередньо зацікавлене.

Все це вже саме по собі, навіть на базі вже існуючих технологій, зменшить необхідний суспільний робочий час і створить простір для існування вільного часу для справжнього людського життя, тобто власне для комунізму.

Рух у напрямку до комунізму вказує ще одна важлива тенденція, яка дозріває в надрах капіталізму і гальмується ним лише тому, що капітал є домінуючою соціальною формою. «Велика історична роль капіталу полягає у створенні цього додаткового праці, надлишкового з точки зору лише споживної вартості, з точки зору простого підтримання існування робітника, і історичне призначення капіталу буде виконано тоді, коли, з одного боку, потреби будуть розвинені настільки, що сам додатковий працю, праця за межами абсолютно необхідної для життя, стане загальною потребою, що випливає з самих індивідуальних потреб людей, і коли, з іншого боку, загальна працьовитість завдяки суворій дисципліні капіталу, через яку пройшли наступні один за одним покоління, розвинеться як загальне надбання нового покоління, – коли, нарешті, ця загальна працьовитість, завдяки розвитку продуктивних сил праці, які постійно підштовхуються капіталом, одержимим безмежною пристрастю до збагачення і діючим в таких умовах, в яких він лише й може реалізувати цю пристрасть, призведе до того, що, з одного боку, володіння загальним багатством і збереження його вимагатимуть від усього суспільства лише порівняно незначної кількості робочого часу і що, з іншого боку, суспільство, що працює, буде науково ставитися до процесу свого прогресуючого відтворення, свого відтворення у все більшій кількості; – отже, тоді, коли припиниться така праця, при якій людина сама робить те, що вона може змусити речі робити для себе, для людини», – писав Маркс. Ми ж додамо, що це стало б завершенням усього періоду існування та розвитку суспільного поділу праці.

Ідея автоматизованої фабрики, що виникла ще за часів Маркса, та ідея широкого використання можливостей природи за допомогою знань про неї на благо всіх людей охоплює реально існуючі тенденції до створення умов для комунізму. Комунізм настане в міру того, як речі все більше і більше виконуватимуть роботу, яку вони можуть виконати для людини і без участі людини, все менше і менше буде вільного людського праці, необхідного для відтворення людського життя, і все більше управління людьми перетворюватиметься на управління речами. Кожна людина зможе все більше і більше часу присвячувати вільній грі фізичних і розумових сил, або, іншими словами, тому, щоб займатися тим, що їй «до душі», на користь усього суспільства. Останнє і є сутністю комуністичної діяльності.

Історична роль капіталу як важливої суспільної форми організації для розвитку продуктивних сил не може бути реалізована в умовах капіталізму. За повної та безперервної влади капітал постійно відтворює матеріальний характер праці, знову й знову гальмуючи розвиток продуктивних сил там, де вимальовується тенденція до їх подолання. Тому в умовах капіталізму наука не може остаточно перетворитися на безпосередню продуктивну силу, а відтворюючи матеріальний характер праці, відтворюються, зокрема, і найбільш відсталі його форми. Повне знищення капіталу як суспільної форми організації – найважливіше питання перехідного періоду до комуністичного суспільства. У цей період необхідно свідомо доводити рух капіталу до передачі чисто матеріального виробництва під управління машин, до виведення людської діяльності за межі матеріального характеру праці.

При цьому не слід забувати про те, що капітал водночас відтворює й протилежну тенденцію, яка веде суспільство назад до капіталізму. І саме з нею, хоча вона нерозривно пов’язана з першою, соціалістичне суспільство може і повинно відкрито боротися, щоб перехід до комунізму взагалі міг відбутися. Лінія фронту між старим капіталістичним і новим комуністичним суспільством об’єктивно проходить там, де є всі матеріальні передумови для того, щоб матеріальну працю раз і назавжди передати машинам, і там, де цього поки ще не можна зробити, не довівши працю до крайнього ступеня поділу та створення, таким чином, матеріальних передумов, щоб передати її машинам. Там, де можлива повна автоматизація чисто матеріального виробництва речей, капітал як суспільне відношення та відповідний характер праці втрачає свою актуальність і тому об’єктивно може і повинен бути знищений. І доки це не буде здійснено, капітал, як істотна суспільна форма відносин, все одно відтворюватиметься.

Відоме питання антикомуністично налаштованого пересічного громадянина: «Хто митиме туалети, якщо кожен займатиметься улюбленою справою?» – втрачає будь-який сенс. Його пафос знецінюється не лише тому, що існує величезна різниця між тим, щоб час від часу відволікатися від улюбленої справи для виконання цього неприємного, але необхідного суспільного обов’язку, і тим, щоб бути на все життя прикутим до нього, але й тому, що зараз вже існують туалети, які самі себе миють. І лише панування приватної власності та те, що робоча сила прибиральників, тобто час життя людей, виявляється в цих умовах дешевшим за речі, заважає їхньому повсюдному поширенню.

За комунізму ж людина, тобто час і наповненість її життя, стає найвищою цінністю, а всі речі мають значення лише в тій мірі, в якій вони забезпечують справжнє людське життя кожного члена суспільства. Сам факт того, що жива людина повинна витрачати час свого життя на нетворчу працю, яку можна передати машинам, є нестерпним для суспільства, де людина є найвищою цінністю. Такий стан речей має бути знищено, і лише після цього можна буде говорити про комунізм, розвинений на власній основі.

Проте слід визнати, що в перехідний період до комунізму певний час необхідна, але не зовсім вільна праця все ще існуватиме, поки не буде розроблено відповідні технології для її усунення та не буде забезпешено їхнє повсюдне впровадження. Але усунення панування капіталу дозволить обмежити робочий день необхідною працею, виключивши працю, яка має сенс лише в рамках руху вартості до самозростання, результати якої постійно знищуються під час криз або просто як «витрати». На перших етапах необхідна (у сенсі ще не повністю вільної) праця розширить свої межі. Це необхідно для забезпечення суспільних фондів та умов для подальшого розвитку. Але загальний напрямок руху суспільства до комунізму, що включає відповідний розвиток продуктивних сил, передбачає зведення такої праці до все менших і менших обсягів. За певних меж таку працю взагалі можна не враховувати, оскільки вона вже не має жодного суттєвого впливу на загальний хід людського життя. Найважливіше в необхідній праці під час переходу до комунізму, про яку йшлося вище, полягає в тому, що ця, ще не зовсім вільна праця, є працею, спрямованою на створення вільного часу та формування особистості, здатної вільно діяти в суспільному масштабі.
\end{document}
